% --- Archon physics formalization source begin ---
% archon:physics
% archon:covers PhyXMiniProblems/problem_phyx_mini_0225.lean
% archon:source-report reports/phyx_mini/problem_phyx_mini_0225.source.json

\chapter{Physics problem phyx\_mini\_0225}
\label{ch:PhyXMiniProblems_problem_phyx_mini_0225}

\paragraph{Problem source.}
\#\# Physical scenario

While driving behind a car traveling at \textbackslash{}( 3.00\textbackslash{},\textbackslash{}mathrm\{m/s\} \textbackslash{}), you notice that one of the car's tires has a small hemispherical bump on its rim as shown in figure. The radii of the car's tires are \textbackslash{}( 0.300\textbackslash{},\textbackslash{}mathrm\{m\} \textbackslash{}).

\#\# Question

What is the bump's period of oscillation?

\#\# Answer choices

- A: A: 0.928s

- B: B: 0.828s

- C: C: 0.728s

- D: D: 0.628s

\#\# Figure caption (auxiliary; use the image as primary evidence)

The image depicts the side view of a tire with a distinct tread pattern consisting of interconnected zigzag lines creating a chevron shape. On the left side of the tire, there's a protrusion labeled with an arrow pointing to it, and the word "Bump" is used to describe this feature. The tire is positioned on a flat surface, indicating its contact with the ground. The overall tone of the image is grayscale, with the tread pattern being a darker shade than the bump.

\#\# Dataset metadata

Wave Properties; Physical Model Grounding Reasoning, Multi-Formula Reasoning

\paragraph{Recorded answer/context.}
D: D: 0.628s

\paragraph{Figure/image path.}
/root/proposal\_for\_physic/science-mango/phyx\_mini\_run/phyx\_data/test\_image/225.png

\paragraph{Formalization target.}
create a compiling Lean file with sorry bodies at `PhyXMiniProblems/problem\_phyx\_mini\_0225.lean`. The Lean declarations must preserve the physical quantities, dimensions or dimensional roles, figure labels, governing-law hypotheses, and final relation expressed by this problem.
Use Mathlib/Physlib names found through LeanExplore where available. If a physics API is missing, introduce faithful local abstractions rather than scalar placeholder aliases.

\begin{theorem}[Physics formalization target]
\label{thm:physics:phyx_mini_0225:target}
\lean{PhyXMiniProblems.ProblemPhyXMini0225.problem_phyx_mini_0225}
\uses{def:physics:phyx-mini-0225:phyxminiproblems-problemphyxmini0225-durationinseconds, def:physics:phyx-mini-0225:phyxminiproblems-problemphyxmini0225-rollingtirebumpsetup, def:physics:phyx-mini-0225:phyxminiproblems-problemphyxmini0225-matchesproblemdescription, def:physics:phyx-mini-0225:phyxminiproblems-problemphyxmini0225-matchesprimaryfigure, def:physics:phyx-mini-0225:phyxminiproblems-problemphyxmini0225-hasphysicalrollingparameters, def:physics:phyx-mini-0225:phyxminiproblems-problemphyxmini0225-satisfiesrollingtirekinematics, def:physics:phyx-mini-0225:phyxminiproblems-problemphyxmini0225-recordedanswerchoice, def:physics:phyx-mini-0225:phyxminiproblems-problemphyxmini0225-matchesdisplayedprecision, def:physics:phyx-mini-0225:phyxminiproblems-problemphyxmini0225-isuniqueclosestanswer}
The assigned autoformalize agent should translate this physics problem into Lean declarations in the covered file, with theorem and lemma proof bodies written as `by sorry`.
\end{theorem}
\begin{proof}
This is an autoformalization task, not a proof task. Produce statements that can later be proved by the physics prover without weakening the physical model.
\end{proof}

% --- Archon declaration topology begin ---
\section{Declaration topology}
\begin{definition}[Length Quantity]
  \label{def:physics:phyx-mini-0225:phyxminiproblems-problemphyxmini0225-lengthquantity}
  \lean{PhyXMiniProblems.ProblemPhyXMini0225.LengthQuantity}
  A nonnegative physical length, used for the tire radius
\end{definition}
\begin{proof}
  This is the declared model definition of Length Quantity.
\end{proof}
\begin{definition}[Duration Quantity]
  \label{def:physics:phyx-mini-0225:phyxminiproblems-problemphyxmini0225-durationquantity}
  \lean{PhyXMiniProblems.ProblemPhyXMini0225.DurationQuantity}
  A nonnegative physical duration, used for the bump's oscillation period
\end{definition}
\begin{proof}
  This is the declared model definition of Duration Quantity.
\end{proof}
\begin{definition}[Angular Speed Quantity]
  \label{def:physics:phyx-mini-0225:phyxminiproblems-problemphyxmini0225-angularspeedquantity}
  \lean{PhyXMiniProblems.ProblemPhyXMini0225.AngularSpeedQuantity}
  The magnitude of a physical angular speed. Radians are dimensionless, so its physical dimension is inverse time
\end{definition}
\begin{proof}
  This is the declared model definition of Angular Speed Quantity.
\end{proof}
\begin{definition}[Speed Quantity]
  \label{def:physics:phyx-mini-0225:phyxminiproblems-problemphyxmini0225-speedquantity}
  \lean{PhyXMiniProblems.ProblemPhyXMini0225.SpeedQuantity}
  A nonnegative physical translational speed
\end{definition}
\begin{proof}
  This is the declared model definition of Speed Quantity.
\end{proof}
\begin{definition}[length Readout]
  \label{def:physics:phyx-mini-0225:phyxminiproblems-problemphyxmini0225-lengthreadout}
  \lean{PhyXMiniProblems.ProblemPhyXMini0225.lengthReadout}
  \uses{def:physics:phyx-mini-0225:phyxminiproblems-problemphyxmini0225-lengthquantity}
  Scalar readout of a physical length in a coherent unit system
\end{definition}
\begin{proof}
  This is the declared model definition of length Readout.
\end{proof}
\begin{definition}[duration Readout]
  \label{def:physics:phyx-mini-0225:phyxminiproblems-problemphyxmini0225-durationreadout}
  \lean{PhyXMiniProblems.ProblemPhyXMini0225.durationReadout}
  \uses{def:physics:phyx-mini-0225:phyxminiproblems-problemphyxmini0225-durationquantity}
  Scalar readout of a physical duration in a coherent unit system
\end{definition}
\begin{proof}
  This is the declared model definition of duration Readout.
\end{proof}
\begin{definition}[angular Speed Readout]
  \label{def:physics:phyx-mini-0225:phyxminiproblems-problemphyxmini0225-angularspeedreadout}
  \lean{PhyXMiniProblems.ProblemPhyXMini0225.angularSpeedReadout}
  \uses{def:physics:phyx-mini-0225:phyxminiproblems-problemphyxmini0225-angularspeedquantity}
  Scalar readout of an angular-speed magnitude in radians per time unit
\end{definition}
\begin{proof}
  This is the declared model definition of angular Speed Readout.
\end{proof}
\begin{definition}[speed Readout]
  \label{def:physics:phyx-mini-0225:phyxminiproblems-problemphyxmini0225-speedreadout}
  \lean{PhyXMiniProblems.ProblemPhyXMini0225.speedReadout}
  \uses{def:physics:phyx-mini-0225:phyxminiproblems-problemphyxmini0225-speedquantity}
  Scalar readout of a translational speed in a coherent unit system
\end{definition}
\begin{proof}
  This is the declared model definition of speed Readout.
\end{proof}
\begin{definition}[length In Meters]
  \label{def:physics:phyx-mini-0225:phyxminiproblems-problemphyxmini0225-lengthinmeters}
  \lean{PhyXMiniProblems.ProblemPhyXMini0225.lengthInMeters}
  \uses{def:physics:phyx-mini-0225:phyxminiproblems-problemphyxmini0225-lengthquantity, def:physics:phyx-mini-0225:phyxminiproblems-problemphyxmini0225-lengthreadout}
  Read a tire radius in metres
\end{definition}
\begin{proof}
  This is the declared model definition of length In Meters.
\end{proof}
\begin{definition}[duration In Seconds]
  \label{def:physics:phyx-mini-0225:phyxminiproblems-problemphyxmini0225-durationinseconds}
  \lean{PhyXMiniProblems.ProblemPhyXMini0225.durationInSeconds}
  \uses{def:physics:phyx-mini-0225:phyxminiproblems-problemphyxmini0225-durationquantity, def:physics:phyx-mini-0225:phyxminiproblems-problemphyxmini0225-durationreadout}
  Read a duration in seconds
\end{definition}
\begin{proof}
  This is the declared model definition of duration In Seconds.
\end{proof}
\begin{definition}[angular Speed In Radians Per Second]
  \label{def:physics:phyx-mini-0225:phyxminiproblems-problemphyxmini0225-angularspeedinradianspersecond}
  \lean{PhyXMiniProblems.ProblemPhyXMini0225.angularSpeedInRadiansPerSecond}
  \uses{def:physics:phyx-mini-0225:phyxminiproblems-problemphyxmini0225-angularspeedquantity, def:physics:phyx-mini-0225:phyxminiproblems-problemphyxmini0225-angularspeedreadout}
  Read an angular speed in radians per second
\end{definition}
\begin{proof}
  This is the declared model definition of angular Speed In Radians Per Second.
\end{proof}
\begin{definition}[speed In Meters Per Second]
  \label{def:physics:phyx-mini-0225:phyxminiproblems-problemphyxmini0225-speedinmeterspersecond}
  \lean{PhyXMiniProblems.ProblemPhyXMini0225.speedInMetersPerSecond}
  \uses{def:physics:phyx-mini-0225:phyxminiproblems-problemphyxmini0225-speedquantity, def:physics:phyx-mini-0225:phyxminiproblems-problemphyxmini0225-speedreadout}
  Read a translational speed in metres per second
\end{definition}
\begin{proof}
  This is the declared model definition of speed In Meters Per Second.
\end{proof}
\begin{definition}[Figure Label]
  \label{def:physics:phyx-mini-0225:phyxminiproblems-problemphyxmini0225-figurelabel}
  \lean{PhyXMiniProblems.ProblemPhyXMini0225.FigureLabel}
  The label printed beside the protrusion in the supplied image
\end{definition}
\begin{proof}
  This is the declared model definition of Figure Label.
\end{proof}
\begin{definition}[Figure Element]
  \label{def:physics:phyx-mini-0225:phyxminiproblems-problemphyxmini0225-figureelement}
  \lean{PhyXMiniProblems.ProblemPhyXMini0225.FigureElement}
  Distinct physical elements visible in the supplied image
\end{definition}
\begin{proof}
  This is the declared model definition of Figure Element.
\end{proof}
\begin{definition}[Figure Region]
  \label{def:physics:phyx-mini-0225:phyxminiproblems-problemphyxmini0225-figureregion}
  \lean{PhyXMiniProblems.ProblemPhyXMini0225.FigureRegion}
  Qualitative image regions used to locate the marked protrusion
\end{definition}
\begin{proof}
  This is the declared model definition of Figure Region.
\end{proof}
\begin{definition}[Bump Shape]
  \label{def:physics:phyx-mini-0225:phyxminiproblems-problemphyxmini0225-bumpshape}
  \lean{PhyXMiniProblems.ProblemPhyXMini0225.BumpShape}
  Shape assigned to the bump by the problem statement
\end{definition}
\begin{proof}
  This is the declared model definition of Bump Shape.
\end{proof}
\begin{definition}[Relative Feature Size]
  \label{def:physics:phyx-mini-0225:phyxminiproblems-problemphyxmini0225-relativefeaturesize}
  \lean{PhyXMiniProblems.ProblemPhyXMini0225.RelativeFeatureSize}
  Relative size assigned to the bump by the problem statement
\end{definition}
\begin{proof}
  This is the declared model definition of Relative Feature Size.
\end{proof}
\begin{definition}[Tread Pattern]
  \label{def:physics:phyx-mini-0225:phyxminiproblems-problemphyxmini0225-treadpattern}
  \lean{PhyXMiniProblems.ProblemPhyXMini0225.TreadPattern}
  Qualitative tread geometry visible in the primary image
\end{definition}
\begin{proof}
  This is the declared model definition of Tread Pattern.
\end{proof}
\begin{definition}[Tire Orientation]
  \label{def:physics:phyx-mini-0225:phyxminiproblems-problemphyxmini0225-tireorientation}
  \lean{PhyXMiniProblems.ProblemPhyXMini0225.TireOrientation}
  Orientation of the tire with respect to the road
\end{definition}
\begin{proof}
  This is the declared model definition of Tire Orientation.
\end{proof}
\begin{definition}[Ground Shape]
  \label{def:physics:phyx-mini-0225:phyxminiproblems-problemphyxmini0225-groundshape}
  \lean{PhyXMiniProblems.ProblemPhyXMini0225.GroundShape}
  Shape of the supporting road surface
\end{definition}
\begin{proof}
  This is the declared model definition of Ground Shape.
\end{proof}
\begin{definition}[Observation Context]
  \label{def:physics:phyx-mini-0225:phyxminiproblems-problemphyxmini0225-observationcontext}
  \lean{PhyXMiniProblems.ProblemPhyXMini0225.ObservationContext}
  The observer's relation to the moving car in the scenario
\end{definition}
\begin{proof}
  This is the declared model definition of Observation Context.
\end{proof}
\begin{definition}[Bump Observable]
  \label{def:physics:phyx-mini-0225:phyxminiproblems-problemphyxmini0225-bumpobservable}
  \lean{PhyXMiniProblems.ProblemPhyXMini0225.BumpObservable}
  The periodic observable whose period the question requests
\end{definition}
\begin{proof}
  This is the declared model definition of Bump Observable.
\end{proof}
\begin{definition}[Tire Figure]
  \label{def:physics:phyx-mini-0225:phyxminiproblems-problemphyxmini0225-tirefigure}
  \lean{PhyXMiniProblems.ProblemPhyXMini0225.TireFigure}
  \uses{def:physics:phyx-mini-0225:phyxminiproblems-problemphyxmini0225-figurelabel, def:physics:phyx-mini-0225:phyxminiproblems-problemphyxmini0225-figureelement, def:physics:phyx-mini-0225:phyxminiproblems-problemphyxmini0225-figureregion, def:physics:phyx-mini-0225:phyxminiproblems-problemphyxmini0225-treadpattern, def:physics:phyx-mini-0225:phyxminiproblems-problemphyxmini0225-tireorientation, def:physics:phyx-mini-0225:phyxminiproblems-problemphyxmini0225-groundshape}
  Categorical information read from the supplied tire image
\end{definition}
\begin{proof}
  This is the declared model definition of Tire Figure.
\end{proof}
\begin{definition}[Rolling Tire Bump Setup]
  \label{def:physics:phyx-mini-0225:phyxminiproblems-problemphyxmini0225-rollingtirebumpsetup}
  \lean{PhyXMiniProblems.ProblemPhyXMini0225.RollingTireBumpSetup}
  \uses{def:physics:phyx-mini-0225:phyxminiproblems-problemphyxmini0225-lengthquantity, def:physics:phyx-mini-0225:phyxminiproblems-problemphyxmini0225-durationquantity, def:physics:phyx-mini-0225:phyxminiproblems-problemphyxmini0225-angularspeedquantity, def:physics:phyx-mini-0225:phyxminiproblems-problemphyxmini0225-speedquantity, def:physics:phyx-mini-0225:phyxminiproblems-problemphyxmini0225-bumpshape, def:physics:phyx-mini-0225:phyxminiproblems-problemphyxmini0225-relativefeaturesize, def:physics:phyx-mini-0225:phyxminiproblems-problemphyxmini0225-observationcontext, def:physics:phyx-mini-0225:phyxminiproblems-problemphyxmini0225-bumpobservable, def:physics:phyx-mini-0225:phyxminiproblems-problemphyxmini0225-tirefigure}
  The physical quantities of the rolling-tire observation. In particular, bumpOscillationPeriod is an unconstrained physical duration here; neither its exact value nor any displayed answer is built into the setup
\end{definition}
\begin{proof}
  This is the declared model definition of Rolling Tire Bump Setup.
\end{proof}
\begin{definition}[Matches Problem Description]
  \label{def:physics:phyx-mini-0225:phyxminiproblems-problemphyxmini0225-matchesproblemdescription}
  \lean{PhyXMiniProblems.ProblemPhyXMini0225.MatchesProblemDescription}
  \uses{def:physics:phyx-mini-0225:phyxminiproblems-problemphyxmini0225-lengthinmeters, def:physics:phyx-mini-0225:phyxminiproblems-problemphyxmini0225-speedinmeterspersecond, def:physics:phyx-mini-0225:phyxminiproblems-problemphyxmini0225-rollingtirebumpsetup}
  Problem-text data: the observer follows the car, the tire radius is 0.300 m, the car speed is 3.00 m/s, and the requested observable is the oscillation period of the stated small hemispherical bump. No period value occurs here
\end{definition}
\begin{proof}
  This is the declared model definition of Matches Problem Description.
\end{proof}
\begin{definition}[Matches Primary Figure]
  \label{def:physics:phyx-mini-0225:phyxminiproblems-problemphyxmini0225-matchesprimaryfigure}
  \lean{PhyXMiniProblems.ProblemPhyXMini0225.MatchesPrimaryFigure}
  \uses{def:physics:phyx-mini-0225:phyxminiproblems-problemphyxmini0225-rollingtirebumpsetup}
  Primary-image evidence: the text label points to a bump on the upper-left rim, the chevron-tread tire stands upright, and it contacts a flat ground surface. This predicate contains no kinematic or period conclusion
\end{definition}
\begin{proof}
  This is the declared model definition of Matches Primary Figure.
\end{proof}
\begin{definition}[Has Physical Rolling Parameters]
  \label{def:physics:phyx-mini-0225:phyxminiproblems-problemphyxmini0225-hasphysicalrollingparameters}
  \lean{PhyXMiniProblems.ProblemPhyXMini0225.HasPhysicalRollingParameters}
  \uses{def:physics:phyx-mini-0225:phyxminiproblems-problemphyxmini0225-lengthinmeters, def:physics:phyx-mini-0225:phyxminiproblems-problemphyxmini0225-durationinseconds, def:physics:phyx-mini-0225:phyxminiproblems-problemphyxmini0225-angularspeedinradianspersecond, def:physics:phyx-mini-0225:phyxminiproblems-problemphyxmini0225-speedinmeterspersecond, def:physics:phyx-mini-0225:phyxminiproblems-problemphyxmini0225-rollingtirebumpsetup}
  Physical admissibility for the quantities used in division and for the selected positive rotation branch. These conditions do not assign the requested period a numerical value
\end{definition}
\begin{proof}
  This is the declared model definition of Has Physical Rolling Parameters.
\end{proof}
\begin{definition}[Satisfies Rolling Tire Kinematics]
  \label{def:physics:phyx-mini-0225:phyxminiproblems-problemphyxmini0225-satisfiesrollingtirekinematics}
  \lean{PhyXMiniProblems.ProblemPhyXMini0225.SatisfiesRollingTireKinematics}
  \uses{def:physics:phyx-mini-0225:phyxminiproblems-problemphyxmini0225-lengthreadout, def:physics:phyx-mini-0225:phyxminiproblems-problemphyxmini0225-durationreadout, def:physics:phyx-mini-0225:phyxminiproblems-problemphyxmini0225-angularspeedreadout, def:physics:phyx-mini-0225:phyxminiproblems-problemphyxmini0225-speedreadout, def:physics:phyx-mini-0225:phyxminiproblems-problemphyxmini0225-rollingtirebumpsetup}
  The car carries the wheel center at the vehicle speed; rolling without slipping gives v = ω R; and the marked bump completes one vertical-position cycle per full revolution, giving ω T = 2π. The laws are required in every coherent unit system and contain no numerical period or answer choice
\end{definition}
\begin{proof}
  This is the declared model definition of Satisfies Rolling Tire Kinematics.
\end{proof}
\begin{definition}[Answer Choice]
  \label{def:physics:phyx-mini-0225:phyxminiproblems-problemphyxmini0225-answerchoice}
  \lean{PhyXMiniProblems.ProblemPhyXMini0225.AnswerChoice}
  Labels of the four answers displayed with the problem
\end{definition}
\begin{proof}
  This is the declared model definition of Answer Choice.
\end{proof}
\begin{definition}[seconds]
  \label{def:physics:phyx-mini-0225:phyxminiproblems-problemphyxmini0225-answerchoice-seconds}
  \lean{PhyXMiniProblems.ProblemPhyXMini0225.AnswerChoice.seconds}
  \uses{def:physics:phyx-mini-0225:phyxminiproblems-problemphyxmini0225-answerchoice}
  Seconds printed beside each displayed answer label
\end{definition}
\begin{proof}
  This is the declared model definition of seconds.
\end{proof}
\begin{definition}[recorded Answer Choice]
  \label{def:physics:phyx-mini-0225:phyxminiproblems-problemphyxmini0225-recordedanswerchoice}
  \lean{PhyXMiniProblems.ProblemPhyXMini0225.recordedAnswerChoice}
  \uses{def:physics:phyx-mini-0225:phyxminiproblems-problemphyxmini0225-answerchoice}
  Dataset metadata: the recorded answer is choice D; this is never a premise
\end{definition}
\begin{proof}
  This is the declared model definition of recorded Answer Choice.
\end{proof}
\begin{definition}[Matches Displayed Precision]
  \label{def:physics:phyx-mini-0225:phyxminiproblems-problemphyxmini0225-matchesdisplayedprecision}
  \lean{PhyXMiniProblems.ProblemPhyXMini0225.MatchesDisplayedPrecision}
  \uses{def:physics:phyx-mini-0225:phyxminiproblems-problemphyxmini0225-durationquantity, def:physics:phyx-mini-0225:phyxminiproblems-problemphyxmini0225-durationinseconds, def:physics:phyx-mini-0225:phyxminiproblems-problemphyxmini0225-answerchoice}
  A duration agrees with a value displayed to the nearest millisecond. The half-millisecond tolerance models the precision of the answer list
\end{definition}
\begin{proof}
  This is the declared model definition of Matches Displayed Precision.
\end{proof}
\begin{definition}[Is Unique Closest Answer]
  \label{def:physics:phyx-mini-0225:phyxminiproblems-problemphyxmini0225-isuniqueclosestanswer}
  \lean{PhyXMiniProblems.ProblemPhyXMini0225.IsUniqueClosestAnswer}
  \uses{def:physics:phyx-mini-0225:phyxminiproblems-problemphyxmini0225-durationquantity, def:physics:phyx-mini-0225:phyxminiproblems-problemphyxmini0225-durationinseconds, def:physics:phyx-mini-0225:phyxminiproblems-problemphyxmini0225-answerchoice}
  A displayed choice is uniquely closest to the model's physical period
\end{definition}
\begin{proof}
  This is the declared model definition of Is Unique Closest Answer.
\end{proof}
\begin{lemma}[bump Oscillation Period formula]
  \label{lem:physics:phyx-mini-0225:phyxminiproblems-problemphyxmini0225-bumposcillationperiod-formula}
  \lean{PhyXMiniProblems.ProblemPhyXMini0225.bumpOscillationPeriod_formula}
  \uses{def:physics:phyx-mini-0225:phyxminiproblems-problemphyxmini0225-lengthinmeters, def:physics:phyx-mini-0225:phyxminiproblems-problemphyxmini0225-durationinseconds, def:physics:phyx-mini-0225:phyxminiproblems-problemphyxmini0225-speedinmeterspersecond, def:physics:phyx-mini-0225:phyxminiproblems-problemphyxmini0225-rollingtirebumpsetup, def:physics:phyx-mini-0225:phyxminiproblems-problemphyxmini0225-hasphysicalrollingparameters, def:physics:phyx-mini-0225:phyxminiproblems-problemphyxmini0225-satisfiesrollingtirekinematics}
  For any physically admissible rolling setup satisfying the two kinematic laws, eliminating angular speed gives T = 2πR/v in SI readouts
\end{lemma}
\begin{proof}
  The conclusion follows from the governing relations and figure readouts named in the statement of bump Oscillation Period formula.
\end{proof}
\begin{lemma}[bump Oscillation Period exact]
  \label{lem:physics:phyx-mini-0225:phyxminiproblems-problemphyxmini0225-bumposcillationperiod-exact}
  \lean{PhyXMiniProblems.ProblemPhyXMini0225.bumpOscillationPeriod_exact}
  \uses{def:physics:phyx-mini-0225:phyxminiproblems-problemphyxmini0225-durationinseconds, def:physics:phyx-mini-0225:phyxminiproblems-problemphyxmini0225-rollingtirebumpsetup, def:physics:phyx-mini-0225:phyxminiproblems-problemphyxmini0225-matchesproblemdescription, def:physics:phyx-mini-0225:phyxminiproblems-problemphyxmini0225-hasphysicalrollingparameters, def:physics:phyx-mini-0225:phyxminiproblems-problemphyxmini0225-satisfiesrollingtirekinematics}
  Substituting R = 0.300 m and v = 3.00 m/s into the derived rolling-period formula gives the exact ideal-model period π/5 s
\end{lemma}
\begin{proof}
  The conclusion follows from the governing relations and figure readouts named in the statement of bump Oscillation Period exact.
\end{proof}
% --- Archon declaration topology end ---
% --- Archon physics formalization source end ---
